%==============================================================================
% 5. RESULTADOS Y DISCUSION
%   -- Version corregida y adaptada a los resultados obtenidos mediante
%      ajuste por MINIMOS CUADRADOS (Python/NumPy/SciPy), en sustitucion
%      del metodo grafico manual de rectas limite, y con los 3 casos de
%      estudio completos (sin clips, con 2 clips, con 3 clips).
%==============================================================================
\section{Resultados y discusion}\label{sec:resultados}

A partir de los datos de posicion y velocidad obtenidos con Tracker 6.3.5 para
los tres casos de estudio (sin clips, con 2 clips y con 3 clips, cada uno con
tres repeticiones), se generaron las curvas de velocidad respecto al tiempo
y se procesaron en Python 3.11.4. Los resultados del ajuste experimental se
presentan en la seccion~\ref{sec:ajuste_experimental}, mientras que en la
seccion~\ref{sec:analisis_comparativo} se realiza el analisis comparativo
entre estos resultados y los tres modelos teoricos desarrollados en la
seccion~\ref{sec:modelo} (caida libre, arrastre viscoso lineal y arrastre de
Newton).

\subsection{Analisis experimental}\label{sec:ajuste_experimental}

La validacion de los modelos matematicos desarrollados en la
Seccion~\ref{sec:modelo} requiere un analisis riguroso de las incertidumbres
experimentales. Siguiendo los principios establecidos por
Taylor~\cite{taylor1997}, ninguna medicion esta libre de incertidumbre, y la
comparacion entre la teoria y la practica debe realizarse evaluando si los
valores teoricos caen dentro de los intervalos de confianza experimentales.

\subsubsection{Tratamiento estadistico y propagacion de errores}

Al repetir el experimento tres veces bajo las mismas condiciones para cada
caso, se observo una dispersion en los valores de la velocidad debido a
errores aleatorios inherentes al proceso de diferenciacion numerica
($v \approx \Delta x/\Delta t$) realizado por Tracker. Para cuantificar esta
dispersion se utilizo el criterio de la \textbf{Desviacion Absoluta Maxima
(d.a.m.)}~\cite{manualgrafico}. Dado que las tres repeticiones de cada caso
no se muestrearon exactamente en los mismos instantes (variaciones propias
del video a $30\text{ fps}$), las tres series se interpolaron sobre una
malla de tiempo comun $t_i$ antes de promediar. Si $\bar{v}_i$ es el valor
promedio (mejor estimacion) de la velocidad interpolada en el instante
$t_i$, la incertidumbre $\delta v_i$ se definio como:

\begin{equation}
    \delta v_i = \max \left( | \bar{v}_i - v_{i,1} |,\ | \bar{v}_i - v_{i,2} |,\ | \bar{v}_i - v_{i,3} | \right)
\end{equation}

El resultado experimental se reporto en la forma estandar
$v_i = \bar{v}_i \pm \delta v_i$. Es importante notar que el proceso de
derivacion numerica amplifica el ruido de la posicion, resultando en barras
de error ($\delta v_i$) proporcionalmente grandes, lo cual es consistente
con la teoria de propagacion de errores~\cite{taylor1997}.

\subsubsection{Ajuste por minimos cuadrados}

Para un objeto en movimiento uniformemente acelerado, el modelo matematico
predice una relacion lineal de la forma $v(t) = v_0 + at$. En esta version
del analisis, en lugar del metodo grafico manual de rectas limite, la recta
$v(t)=v_0+at$ se ajusto de forma objetiva mediante el \textbf{metodo de
minimos cuadrados}, implementado en Python (\texttt{NumPy}). Para cada
instante $t_i$ dentro del intervalo de validez lineal se minimizo la suma de
los residuos cuadraticos:

\begin{equation}
    S(v_0,a) = \sum_{i=1}^{n}\left[\bar{v}_i - (v_0+at_i)\right]^2
\end{equation}

y la incertidumbre de la pendiente y la ordenada al origen se calcularon con
las expresiones estandar de la regresion lineal~\cite{taylor1997}:

\begin{equation}
    \delta a = s_{\text{res}}\sqrt{\dfrac{n}{n\sum t_i^2-\left(\sum t_i\right)^2}}\, ,
    \qquad
    \delta v_0 = s_{\text{res}}\sqrt{\dfrac{\sum t_i^2}{n\sum t_i^2-\left(\sum t_i\right)^2}}
\end{equation}

donde $s_{\text{res}}$ es la desviacion estandar de los residuos del ajuste.
El intervalo de validez de cada ajuste ($t\le T_{\text{corte}}$) se
determino identificando visualmente el instante en el que los datos
experimentales se desvian de la tendencia lineal, ya sea por el impacto con
el suelo o por la aproximacion a la velocidad terminal.
\\

\textbf{Caso 0, sin clips:}

\begin{figure}[htbp]
  \centering
  \includegraphics[width=0.75\linewidth]{figuras/velocidad_vs_tiempo_sinclips_minimos_cuadrados.png}
  \caption{Velocidad del filtro de café sin clips, con barras de error (d.a.m.) y ajuste por mínimos cuadrados.}
  \label{fig:sin_clips_ajuste}
\end{figure}

El ajuste por minimos cuadrados, restringido al intervalo
$t\in[0,\,1.7]\text{ s}$ ($n=52$ puntos), arrojo:

\begin{equation}
    v(t) = (-0.826 \pm 0.067) + (-0.346 \pm 0.068)\,t \quad \text{[m/s]}, \qquad R^2 = 0.344
\end{equation}

Aplicando las reglas de cifras significativas y redondeo de
Taylor~\cite{taylor1997}, la aceleracion experimental se reporta como:

\begin{equation}
    a_{exp}^{(0)} = -0.3 \pm 0.1 \text{ m/s}^2
\end{equation}

Para tiempos $t>1.7\text{ s}$ los datos experimentales se desvian de la
tendencia lineal (rebote por el impacto con el suelo), por lo que la
ecuacion empirica anterior solo describe la realidad fisica dentro de dicho
intervalo. El valor de $R^2$, relativamente bajo, refleja la gran dispersion
introducida por la diferenciacion numerica en un caso donde la aceleracion
neta es pequena (la resistencia del aire casi compensa el peso incluso para
el filtro solo).
\\

\textbf{Caso 1, con 2 clips:}

\begin{figure}[htbp]
  \centering
  \includegraphics[width=0.75\linewidth]{figuras/velocidad_vs_tiempo_2clips_minimos_cuadrados.png}
  \caption{Velocidad del filtro de café con 2 clips, con barras de error (d.a.m.) y ajuste por mínimos cuadrados.}
  \label{fig:con2_clips_ajuste}
\end{figure}

El ajuste, restringido a $t\in[0,\,1.24]\text{ s}$ ($n=38$ puntos), arrojo:

\begin{equation}
    v(t) = (-1.034 \pm 0.076) + (-0.764 \pm 0.105)\,t \quad \text{[m/s]}, \qquad R^2 = 0.594
\end{equation}

por lo que la aceleracion experimental equivalente se reporta como:

\begin{equation}
    a_{exp}^{(1)} = -0.8 \pm 0.1 \text{ m/s}^2
\end{equation}

Para $t>1.24\text{ s}$ los datos se desvian de la tendencia lineal: la
velocidad comienza a estabilizarse (aproximandose a la velocidad terminal) o
presenta ruido por el impacto con el suelo.
\\

\textbf{Caso 2, con 3 clips:}

\begin{figure}[htbp]
  \centering
  \includegraphics[width=0.75\linewidth]{figuras/velocidad_vs_tiempo_3clips_minimos_cuadrados.png}
  \caption{Velocidad del filtro de café con 3 clips, con barras de error (d.a.m.) y ajuste por mínimos cuadrados.}
  \label{fig:con3_clips_ajuste}
\end{figure}

El ajuste, restringido a $t\in[0,\,1.1]\text{ s}$ ($n=34$ puntos), arrojo:

\begin{equation}
    v(t) = (-0.969 \pm 0.084) + (-1.234 \pm 0.131)\,t \quad \text{[m/s]}, \qquad R^2 = 0.736
\end{equation}

por lo que la aceleracion experimental equivalente se reporta como:

\begin{equation}
    a_{exp}^{(2)} = -1.2 \pm 0.1 \text{ m/s}^2
\end{equation}

En este caso el intervalo de validez lineal es el mas corto de los tres
($t\le 1.1\text{ s}$): al ser el caso con mayor peso, el objeto alcanza su
velocidad terminal (y por tanto se aleja del comportamiento lineal) en menos
tiempo, y tambien llega antes al suelo.
\\

La Tabla~\ref{tab:resumen_ajuste_lineal} resume los parametros obtenidos
para los tres casos. Se observa que, conforme aumenta la masa del sistema
(mas clips), la magnitud de la pendiente $a_{exp}$ crece de forma monotona,
y el coeficiente de determinacion $R^2$ mejora, lo cual es consistente con
una mayor relacion senal/ruido en los casos de caida mas rapida.

\begin{table}[htbp]
  \centering
  \caption{Resumen del ajuste lineal por minimos cuadrados, $v(t)=v_0+at$, para los tres casos de estudio.}
  \label{tab:resumen_ajuste_lineal}
  \begin{tabular}{@{}lccccc@{}}
    \toprule
    Caso & $m$ [g] & $v_0$ [m/s] & $a_{exp}$ [m/s$^2$] & $R^2$ & Intervalo de validez [s] \\
    \midrule
    Sin clips    & 0.93 & $-0.826 \pm 0.067$ & $-0.3 \pm 0.1$ & 0.344 & $[0,\,1.70]$ \\
    Con 2 clips  & 2.01 & $-1.034 \pm 0.076$ & $-0.8 \pm 0.1$ & 0.594 & $[0,\,1.24]$ \\
    Con 3 clips  & 2.55 & $-0.969 \pm 0.084$ & $-1.2 \pm 0.1$ & 0.736 & $[0,\,1.10]$ \\
    \bottomrule
  \end{tabular}
\end{table}

\subsection{Analisis comparativo}\label{sec:analisis_comparativo}

Con los datos promediados $\bar{v}(t)$ obtenidos en la
seccion~\ref{sec:ajuste_experimental}, se realizo un analisis comparativo
entre los resultados experimentales y los tres modelos teoricos
desarrollados en la seccion~\ref{sec:modelo}:

\begin{itemize}
  \item \textbf{Caida libre} (sin resistencia del aire): $v(t) = -gt$, sin parametros libres.
  \item \textbf{Arrastre viscoso lineal} (Caso 2, seccion~\ref{sec:region2}): $v(t) = -\dfrac{g}{p}\left(1-e^{-pt}\right)$, con $p=b/m$ como unico parametro libre.
  \item \textbf{Arrastre de Newton} (Caso 3, seccion~\ref{sec:region3_newton}): $v(t) = -v_t\tanh\!\left(\dfrac{g}{v_t}t\right)$, con la velocidad terminal $v_t$ como unico parametro libre.
\end{itemize}

A diferencia del ajuste de la seccion~\ref{sec:ajuste_experimental}, aqui el
parametro libre de cada modelo ($p$ o $v_t$) no describe una recta sino una
curva no lineal (exponencial o tangente hiperbolica); por ello se ajusto
mediante \textbf{minimos cuadrados no lineales} (rutina
\texttt{scipy.optimize.curve\_fit} de Python), utilizando como referencia la
velocidad promedio experimental $\bar{v}(t)$ de cada caso, ponderada por su
incertidumbre $\delta v_i$. A partir del parametro $p$ ajustado se obtuvo el
coeficiente de arrastre viscoso $b=pm$; a partir de $v_t$ se obtuvo la
cantidad $c\rho A = 2mg/v_t^2$, que agrupa el coeficiente de arrastre, la
densidad del aire y el area frontal del objeto (no separables con la
instrumentacion disponible). Para la masa de cada caso se utilizaron los
valores medidos $m_{filtro}=\cantidad{0.00093}{kg}$ y
$m_{clip}=\cantidad{0.00054}{kg}$ (Tabla~\ref{tab:parametros}).

Los resultados se muestran en la Figura~\ref{fig:comparacion_modelos} y se
resumen en la Tabla~\ref{tab:resumen_modelos}.

\begin{figure}[htbp]
  \centering
  \includegraphics[width=\linewidth]{figuras/comparacion_modelos_teoricos.png}
  \caption{Comparación entre los datos experimentales (promedio ± d.a.m.) y los tres modelos teóricos ajustados: caída libre, arrastre viscoso lineal y arrastre de Newton, para los tres casos de estudio.}
  \label{fig:comparacion_modelos}
\end{figure}

\begin{table}[htbp]
  \centering
  \caption{Coeficiente de determinacion $R^2$ y parametros fisicos ajustados para cada modelo teorico, en los tres casos de estudio.}
  \label{tab:resumen_modelos}
  \begin{tabular}{@{}lcccccc@{}}
    \toprule
     & \multicolumn{1}{c}{Caida libre} & \multicolumn{2}{c}{Arrastre lineal} & \multicolumn{2}{c}{Arrastre de Newton} \\
    \cmidrule(lr){2-2}\cmidrule(lr){3-4}\cmidrule(lr){5-6}
    Caso & $R^2$ & $R^2$ & $b$ [kg/s] & $R^2$ & $v_t$ [m/s] ($c\rho A$ [kg/m]) \\
    \midrule
    Sin clips   & $-856.8$ & 0.690 & $7.77\times10^{-3}$ & 0.446 & $1.08$ $(1.56\times10^{-2})$ \\
    Con 2 clips & $-239.7$ & 0.849 & $1.10\times10^{-2}$ & 0.788 & $1.75$ $(1.29\times10^{-2})$ \\
    Con 3 clips & $-98.2$  & 0.912 & $1.26\times10^{-2}$ & 0.836 & $1.95$ $(1.32\times10^{-2})$ \\
    \bottomrule
  \end{tabular}
\end{table}

\textbf{Caida libre.} En los tres casos el modelo de caida libre se rechaza
categoricamente: $R^2$ es fuertemente negativo (indica que el modelo ajusta
peor que una simple recta horizontal en el promedio de los datos). Esto es
fisicamente razonable: al tratarse de un objeto muy ligero
($m\sim1\text{--}3\text{ g}$) y de area frontal relativamente grande (un
filtro de papel), la resistencia del aire deja de ser despreciable practicamente
desde el instante inicial de la caida, y el sistema se aproxima con rapidez
a una velocidad terminal muy por debajo de la que predice $v(t)=-gt$.

\textbf{Arrastre lineal vs.\ arrastre de Newton.} En los tres casos, el
modelo de \textbf{arrastre lineal} describe mejor los datos experimentales
que el de \textbf{arrastre de Newton} ($R^2$ mayor en los tres), aunque
ambos modelos siguen razonablemente bien la forma general de la curva,
incluyendo el acercamiento asintotico a la velocidad terminal. Esto sugiere
que, en el intervalo de velocidades alcanzado por el experimento
($\sim 1$--$2\text{ m/s}$), el regimen de arrastre del filtro de cafe se
aproxima mas al regimen viscoso/laminar ($F\propto v$) que al turbulento
($F\propto v^2$); sin embargo, dado que ambos $R^2$ son de magnitud
comparable (particularmente en los casos con clips), no es posible
descartar por completo el modelo de arrastre de Newton solo con este
criterio, y un numero de Reynolds intermedio explicaria que ninguno de los
dos modelos puros ajuste perfectamente.

\textbf{Consistencia fisica de los parametros ajustados.} Tanto $b$
(arrastre lineal) como $c\rho A$ (arrastre de Newton) se mantienen del mismo
orden de magnitud en los tres casos (variando aproximadamente un factor de
1.6 y 1.2, respectivamente, pese a que la masa varia por un factor de 2.7).
Esto es consistente con la fisica del problema: ambos parametros dependen
principalmente de la geometria del objeto (area frontal, forma), la cual es
practicamente la misma en los tres casos (el mismo filtro de cafe, solo se
le agregan clips que aportan masa pero modifican muy poco su area frontal
efectiva), y no deberian depender fuertemente de la masa agregada. La
variacion residual observada puede atribuirse a cambios menores en la
orientacion o en el area efectiva del filtro al agregar los clips, asi como
al ruido propio de la diferenciacion numerica discutido en la
seccion~\ref{sec:ajuste_experimental}.

\textbf{Mejora del ajuste con la masa.} Al igual que en el ajuste lineal de
la seccion~\ref{sec:ajuste_experimental}, el $R^2$ de ambos modelos de
arrastre mejora conforme aumenta la masa (mas clips). Esto se explica porque,
al aumentar el peso, la aceleracion neta y las velocidades alcanzadas son
mayores en comparacion con el ruido de medicion (que es aproximadamente del
mismo orden absoluto en los tres casos), mejorando la relacion senal/ruido y,
por tanto, la calidad del ajuste.
